\documentclass[11pt,a4paper]{article}
\usepackage[margin=2.4cm]{geometry}
\usepackage{amsmath,amssymb,amsthm,mathtools}
\usepackage{booktabs,array,tabularx}
\usepackage{graphicx}
\usepackage{enumitem}
\usepackage[hidelinks]{hyperref}
\usepackage{xcolor}
\usepackage{algorithm}
\usepackage{algpseudocode}

\newtheorem{theorem}{Theorem}
\newtheorem{proposition}[theorem]{Proposition}
\theoremstyle{remark}
\newtheorem{remark}{Remark}

\newcommand{\R}{\mathbb{R}}
\newcommand{\K}{\mathcal{K}}
\newcommand{\norm}[1]{\left\lVert #1\right\rVert}
\newcommand{\ip}[2]{#1^{\mathsf T}#2}
\newcommand{\T}{^{\mathsf T}}

\title{Conjugate Gradient Beyond the Quadratic Case\\[0.3em]
\large Literature map and extensions after U18--U19 (MMC007 Optimization, IMCA)}
\author{Gustavo Lozano}
\date{November 2026}

\sloppy
\begin{document}
\maketitle

\paragraph{Companion material.} All paths are relative to the repository root.
\begin{itemize}[leftmargin=1.5em,itemsep=2pt]
\item Unit labels (U5, U7, U15--U21, in particular U18 and U19) refer to the optimization agenda
\path{master/IMCA/20262/courses/MMC007/exercises/agenda_optimizacion.pdf}.
\item The code and experiments (\texttt{linear\_cg.py}, \texttt{P18.py}, \texttt{P20.py}) are in
\path{master/IMCA/20262/courses/MMC007/exercises/src/}.
\end{itemize}


\begin{abstract}
Units U18--U19 of the course agenda covered linear conjugate gradient (CG) for $\min \tfrac12 x\T Ax - b\T x$ with $A\succ 0$. Experiments P18 and P20 confirmed two facts: CG terminates in at most as many steps as there are distinct eigenvalues of $A$, and its rate $2\bigl(\tfrac{\sqrt\kappa-1}{\sqrt\kappa+1}\bigr)^k$ is far better than steepest descent's $\bigl(\tfrac{\kappa-1}{\kappa+1}\bigr)^k$. This note answers \emph{whether CG can be used beyond quadratics}. The answer is yes, in three different senses: (i)~Krylov relatives of CG solve linear systems that are not symmetric positive definite; (ii)~linear CG becomes the \emph{inner solver} of Newton and trust-region methods for general smooth $f$; (iii)~\emph{nonlinear CG} (Fletcher--Reeves, Polak--Ribi\`ere, Hager--Zhang, \dots) applies the CG recurrence directly to general $f$. For each sense it states the key results, shows a small experiment, and gives an annotated reading list for when the professor's articles arrive.
\end{abstract}

\tableofcontents



%======================================================================
\section{Starting point: what P18 and P20 established}
%======================================================================

With $r_k = Ax_k - b = \nabla f(x_k)$, the (preconditioned, $M\succ0$) linear CG implemented in \texttt{linear\_cg.py} is
\begin{equation}\label{eq:lcg}
\begin{aligned}
&\alpha_k = \frac{\ip{r_k}{y_k}}{\ip{p_k}{Ap_k}},\quad
x_{k+1}=x_k+\alpha_k p_k,\quad
r_{k+1}=r_k+\alpha_k Ap_k,\\
&y_{k+1}=M^{-1}r_{k+1},\quad
\beta_{k+1}=\frac{\ip{r_{k+1}}{y_{k+1}}}{\ip{r_k}{y_k}},\quad
p_{k+1}=-y_{k+1}+\beta_{k+1}p_k .
\end{aligned}
\end{equation}
The two experiments verified:
\begin{itemize}[leftmargin=1.5em]
\item \textbf{P18} ($n=2$): CG reaches $x^*$ in $2$ steps. Steepest descent and coordinate descent zig-zag.
\item \textbf{P20} ($A=\operatorname{diag}(1,12,\dots,100)$, $\kappa=100$): CG terminates at $k=10$ because $A$ has $10$ distinct eigenvalues. Before that, it follows (and beats) the bound $2(9/11)^k$. Steepest descent follows $(99/101)^k$.
\end{itemize}
Both facts come from one property: $x_k$ minimizes $\norm{x-x^*}_A$ over $x_0+\K_k(A,r_0)$, so
\[
\frac{\norm{x_k-x^*}_A}{\norm{x_0-x^*}_A}\le \min_{\substack{P\in\mathcal P_k\\P(0)=1}}\ \max_{\lambda\in\sigma(A)}|P(\lambda)|.
\]
Each generalization below keeps some of these ingredients and drops others. Table~\ref{tab:ingredients} shows which, and it is the most useful way to read the literature.

\begin{table}[h]
\centering\small
\begin{tabularx}{\textwidth}{>{\raggedright\arraybackslash}p{4.2cm}X}
\toprule
Ingredient of linear CG & What happens when it is lost \\
\midrule
$A$ symmetric & No short (3-term) recurrence that is also optimal (Faber--Manteuffel theorem~\cite{FaberManteuffel1984}). One must choose between GMRES (optimal, long recurrence) and BiCG/BiCGStab (short, not optimal).\\
$A$ positive definite & $\norm{\cdot}_A$ is not a norm. Use MINRES/SYMMLQ~\cite{PaigeSaunders1975}, or in optimization \emph{detect} negative curvature $p\T Ap\le 0$ and stop (Steihaug--Toint).\\
Constant Hessian ($f$ quadratic) & Conjugacy holds only approximately. $\alpha_k$ needs a line search, and the different formulas for $\beta_k$ stop being equivalent. This is nonlinear CG (Section~\ref{sec:ncg}).\\
Exact arithmetic & Orthogonality is lost, so convergence is delayed, though not destroyed~\cite{MeurantStrakos2006,CarsonLiesenStrakos2024}.\\
\bottomrule
\end{tabularx}
\caption{Each extension of CG drops one ingredient of the linear method.}
\label{tab:ingredients}
\end{table}

%======================================================================
\section{Sense (i): other linear systems (still linear, but not SPD)}
%======================================================================
\begin{itemize}[leftmargin=1.5em]
\item \textbf{Least squares} $\min\norm{Jx-r}^2$: apply CG to $J\T Jx=J\T r$ without forming $J\T J$ (CGLS/CGNR). LSQR and LSMR are the numerically stable versions~\cite{PaigeSaunders1982,FongSaunders2011}. The effective condition number is $\kappa(J)^2$. This is the linear-algebra core of Gauss--Newton (U21).
\item \textbf{Symmetric indefinite}: MINRES minimizes $\norm{r_k}_2$ over the same Krylov space~\cite{PaigeSaunders1975}.
\item \textbf{Nonsymmetric}: GMRES, BiCG, BiCGStab; see Saad's book~\cite{Saad2003}.
\end{itemize}
These are relatives of CG, not CG itself. Optimization uses them mostly through Gauss--Newton and KKT systems.

%======================================================================
\section{Sense (ii): linear CG inside Newton and trust-region methods}
\label{sec:newtoncg}
%======================================================================
This is the most important practical use of linear CG for general $f$. At $x_k$ the Newton step solves
\[
\nabla^2 f(x_k)\,p = -\nabla f(x_k),
\]
which is exactly a quadratic problem of the U18--U19 type with $A=\nabla^2 f(x_k)$ and $b=-\nabla f(x_k)$. CG needs only the product $v\mapsto\nabla^2 f(x_k)v$, which can be computed by automatic differentiation or a finite difference of gradients, so the Hessian never has to be formed.

\paragraph{Inexact (truncated) Newton.} Stop CG when $\norm{\nabla^2 f_k\,p+\nabla f_k}\le\eta_k\norm{\nabla f_k}$. Dembo, Eisenstat and Steihaug~\cite{DemboEisenstatSteihaug1982} proved that $\eta_k\to0$ gives superlinear convergence and $\eta_k=O(\norm{\nabla f_k})$ gives quadratic convergence (\cite[Thm.~7.2]{NocedalWright2006}). If CG meets $p_j\T\nabla^2 f_k\,p_j\le0$, the Hessian is not positive definite, and CG stops (returning $-\nabla f_k$ on the first iteration). This is N\&W Algorithm~7.1 (line-search Newton--CG); see the survey by Nash~\cite{Nash2000}.

\paragraph{Trust region (connects to U15--U17).} Steihaug~\cite{Steihaug1983} and Toint showed that running preconditioned CG on the model $m_k(p)$ from $p=0$ and stopping when (a)~the iterate leaves the ball $\norm{p}\le\Delta_k$ or (b)~negative curvature appears gives a step at least as good as the Cauchy point. That is enough for the global convergence theory of U17. The first CG iterate \emph{is} the Cauchy point, so the dogleg of U16 is a two-segment version of the CG path. GLTR~\cite{GouldLucidiRomaToint1999} continues past the boundary with Lanczos.

\paragraph{Large-scale learning.} ``Hessian-free'' optimization of neural networks~\cite{Martens2010} is truncated Newton--CG with Gauss--Newton curvature products.

\paragraph{Preconditioning.} The matrix $M$ in \texttt{linear\_cg.py} plays the same role here. In practice $M$ is incomplete Cholesky, a diagonal, an L-BFGS matrix, or a randomized low-rank (Nystr\"om) approximation~\cite{Benzi2002,FrangellaTroppUdell2023}.

%======================================================================
\section{Sense (iii): nonlinear CG for general smooth $f$}
\label{sec:ncg}
%======================================================================
\subsection{The algorithm}
Replace $r_k$ by $g_k=\nabla f(x_k)$, the closed-form $\alpha_k$ by a line search, and keep the direction recurrence:

\begin{algorithm}[h]
\caption{Nonlinear CG (Fletcher--Reeves 1964~\cite{FletcherReeves1964} and successors)}
\begin{algorithmic}[1]
\State Given $x_0$; set $g_0=\nabla f(x_0)$, $d_0=-g_0$.
\For{$k=0,1,2,\dots$ until $\norm{g_k}\le\varepsilon$}
  \State Choose $\alpha_k$ by a (strong) Wolfe line search along $d_k$ \Comment{agenda default: $c_1=10^{-4}$, $c_2=0.1$}
  \State $x_{k+1}=x_k+\alpha_k d_k$, \quad $g_{k+1}=\nabla f(x_{k+1})$, \quad $y_k=g_{k+1}-g_k$
  \State $d_{k+1}=-g_{k+1}+\beta_{k+1}d_k$ \quad (restart: $\beta_{k+1}=0$ every $n$ steps or when $d_{k+1}$ is not a descent direction)
\EndFor
\end{algorithmic}
\end{algorithm}

Memory is a few vectors, and cost per iteration is one gradient plus the line search. This is why nonlinear CG is used when $n$ is very large and even L-BFGS memory is too costly.

\subsection{The choice of $\beta_k$}

\begin{table}[h]
\centering\footnotesize
\begin{tabular}{@{}>{\raggedright\arraybackslash}p{2.6cm}lc>{\raggedright\arraybackslash}p{4.6cm}@{}}
\toprule
Name & $\beta_{k+1}$ & Year & Key property\\
\midrule
Hestenes--Stiefel (HS) & $\dfrac{\ip{g_{k+1}}{y_k}}{\ip{d_k}{y_k}}$ & 1952~\cite{HestenesStiefel1952} & satisfies $\ip{d_{k+1}}{y_k}=0$ (conjugacy)\\[1.2ex]
Fletcher--Reeves (FR) & $\dfrac{\norm{g_{k+1}}^2}{\norm{g_k}^2}$ & 1964~\cite{FletcherReeves1964} & global conv.\ with strong Wolfe, $c_2<\tfrac12$~\cite{AlBaali1985}\\[1.2ex]
Polak--Ribi\`ere--Polyak (PRP) & $\dfrac{\ip{g_{k+1}}{y_k}}{\norm{g_k}^2}$ & 1969~\cite{PolakRibiere1969,Polyak1969} & auto-restart after tiny steps; may cycle~\cite{Powell1984}\\[1.2ex]
PR$+$ & $\max\{\beta^{\rm PRP},0\}$ & 1992~\cite{GilbertNocedal1992} & global conv.\ (Gilbert--Nocedal)\\[1.2ex]
Dai--Yuan (DY) & $\dfrac{\norm{g_{k+1}}^2}{\ip{d_k}{y_k}}$ & 1999~\cite{DaiYuan1999} & descent + conv.\ with \emph{standard} Wolfe\\[1.2ex]
Hager--Zhang (HZ) & $\dfrac{1}{\ip{d_k}{y_k}}\Bigl(y_k-2d_k\dfrac{\norm{y_k}^2}{\ip{d_k}{y_k}}\Bigr)\T g_{k+1}$ & 2005~\cite{HagerZhang2005} & $\ip{g_k}{d_k}\le-\tfrac78\norm{g_k}^2$ for any $\alpha_k$\\
\bottomrule
\end{tabular}
\caption{Main choices of $\beta$. Hundreds of hybrids exist; see the surveys~\cite{HagerZhang2006,Dai2011,NarushimaYabe2014,Andrei2020}.}
\label{tab:beta}
\end{table}

\begin{proposition}[All formulas agree on quadratics]
If $f(x)=\tfrac12x\T Ax-b\T x$ with $A\succ0$ and $\alpha_k$ is the exact minimizer along $d_k$, then $\beta^{\rm HS}=\beta^{\rm FR}=\beta^{\rm PRP}=\beta^{\rm DY}=\beta^{\rm HZ}$, and Algorithm~1 coincides with linear CG~\eqref{eq:lcg} with $M=I$.
\end{proposition}
\begin{proof}[Sketch]
Exact line search gives $\ip{g_{k+1}}{d_k}=0$. Linear CG gives $\ip{g_{k+1}}{g_k}=0$ (residuals are mutually orthogonal). Hence
\[
\ip{g_{k+1}}{y_k}=\norm{g_{k+1}}^2,\qquad
\ip{d_k}{y_k}=\ip{d_k}{g_{k+1}}-\ip{d_k}{g_k}=-\ip{(-g_k+\beta_kd_{k-1})}{g_k}=\norm{g_k}^2,
\]
using $\ip{d_{k-1}}{g_k}=0$. So all numerators equal $\norm{g_{k+1}}^2$ and all denominators equal $\norm{g_k}^2$, which is the FR value. The second HZ term vanishes because $\ip{d_k}{g_{k+1}}=0$. Finally $y_k=\alpha_kAd_k$, so $\ip{d_{k+1}}{Ad_k}=0$ is the HS conjugacy condition.
\end{proof}

So the formulas differ only through how they behave \emph{away} from the quadratic model, with inexact line searches and a Hessian that changes from point to point. That is the whole subject of the literature in Table~\ref{tab:beta}.

\subsection{What is known about convergence}
\begin{itemize}[leftmargin=1.5em]
\item \textbf{Tool:} Zoutendijk's condition $\sum_k\cos^2\theta_k\norm{g_k}^2<\infty$ under Wolfe steps (U5). Every global result proves either $\cos\theta_k$ bounded away from $0$ or $\sum\norm{d_k}^{-2}=\infty$.
\item \textbf{FR:} $\liminf\norm{g_k}=0$ with strong Wolfe, $c_2<1/2$~\cite{AlBaali1985}. In practice FR can get stuck making tiny steps with $\cos\theta_k\approx0$ (N\&W §5.2), so it needs restarts.
\item \textbf{PRP:} Powell~\cite{Powell1984} built a nonconvex example where PRP with \emph{exact} line search cycles without approaching a stationary point. The fix $\beta=\max\{\beta^{\rm PRP},0\}$ (PR$+$) is globally convergent~\cite{GilbertNocedal1992}, and in practice PRP/PR$+$ is much better than FR.
\item \textbf{Rate:} with exact line search and restarts every $n$ steps, the rate is $n$-step quadratic, $\norm{x_{k+n}-x^*}=O(\norm{x_k-x^*}^2)$. Without restarts it is generally only linear, never superlinear~(Crowder--Wolfe; N\&W §5.2).
\item \textbf{Recent worst-case analysis.} Das Gupta, Freund, Sun and Taylor~\cite{DasGuptaEtAl2025} used computer-assisted performance estimation (PEP) to obtain the first non-asymptotic bound for FR and an improved one for PRP on $L$-smooth, $\mu$-strongly convex $f$. Their adversarial examples show that in the worst case these methods are \emph{no better than steepest descent with exact line search}. The P20 advantage ($\sqrt\kappa$ vs.\ $\kappa$) therefore does not carry over to nonlinear CG as a worst-case guarantee. It holds only asymptotically, near $x^*$, where $f$ is nearly quadratic.
\end{itemize}

\subsection{Modern implementations and the bridge to L-BFGS (U20)}
\begin{itemize}[leftmargin=1.5em]
\item \textbf{CG\_DESCENT} (Hager--Zhang): the HZ $\beta$ plus an ``approximate Wolfe'' line search that is accurate near the solution~\cite{HagerZhang2005}. HZ comes from Perry/Shanno's \emph{memoryless BFGS}: a nonlinear CG direction is a quasi-Newton direction that uses only the last pair $(s_k,y_k)$.
\item \textbf{Limited-memory CG}~\cite{HagerZhang2013}: on ill-conditioned problems, rounding makes successive gradients lose the orthogonality of the linear case. The method detects this and restores orthogonality with a short L-BFGS-type memory. It sits between nonlinear CG and L-BFGS, which is exactly the topic of U20.
\item \textbf{CGOPT} (Dai--Kou)~\cite{DaiKou2013}: a CG direction chosen closest to the scaled memoryless BFGS direction. It is competitive with CG\_DESCENT.
\item \textbf{Preconditioning} carries over directly:
\[
d_{k+1}=-M^{-1}g_{k+1}+\beta_{k+1}d_k,\qquad \beta^{\rm PR}_{k+1}=\frac{\ip{g_{k+1}}{M^{-1}y_k}}{\ip{g_k}{M^{-1}g_k}}.
\] This is the same $y=M^{-1}r$ step as in \texttt{linear\_cg.py}.
\end{itemize}

\subsection{Further generalizations (research frontier)}
\begin{itemize}[leftmargin=1.5em]
\item \textbf{Riemannian CG.} For $f$ defined on a manifold (Stiefel, Grassmann, SPD matrices), $d_k$ must be moved to the new tangent space by a vector transport before forming $-\operatorname{grad}f(x_{k+1})+\beta d_k$. Global convergence depends on how the transport is chosen~\cite{RingWirth2012,SatoIwai2015,Sato2022}.
\item \textbf{Hilbert spaces and nonlinear operator equations:} Daniel~\cite{Daniel1967}.
\item \textbf{Monotone nonlinear equations $F(x)=0$:} CG directions combined with hyperplane projection (many recent HZ-type papers).
\item \textbf{High-performance linear CG:} pipelined and communication-avoiding variants hide the global reductions $\ip{r}{y}$ and $\ip{p}{w}$~\cite{GhyselsVanroose2014}. These are relevant if you write your own C++ kernels.
\end{itemize}

%======================================================================
\section{Small experiment}
%======================================================================
\textbf{Setup.} Chained Rosenbrock $f(x)=\sum_{i=1}^{n-1}100(x_{i+1}-x_i^2)^2+(1-x_i)^2$ with $n=100$ and $x_0=(-1.2,1,-1.2,1,\dots)$. Strong Wolfe line search with $c_1=10^{-4}$, $c_2=0.1$ (the agenda default for nonlinear CG). Stop at $\norm{\nabla f}\le10^{-6}$ or after $5000$ iterations. Newton--CG and L-BFGS use the SciPy implementations, which have their own line searches, so their evaluation counts are only indicative.

\begin{figure}[h]
\centering
\includegraphics[width=\textwidth]{ncg_rosen100.pdf}
\caption{Gradient norm on chained Rosenbrock ($n=100$). Every nonlinear CG variant converges; steepest descent stalls at $\norm{\nabla f}\approx3$ after 5000 iterations.}
\label{fig:rosen100}
\end{figure}

\begin{table}[h]
\centering\small
\begin{tabular}{@{}lrrl@{}}
\toprule
Method & Iterations & $f$/$\nabla f$ evals & Comment\\
\midrule
Steepest descent & $>5000$ & 53\,792 & not converged ($\norm{\nabla f}\approx3.0$)\\
FR & 1226 & 14\,431 & long plateau\\
FR, restart every $n$ & 983 & 11\,684 & restarts help FR\\
DY & 1476 & 17\,712 & \\
HS & 872 & 9\,325 & \\
HZ & 906 & 9\,520 & \\
PR$+$ & \textbf{761} & \textbf{8\,160} & best nonlinear CG here\\
L-BFGS ($m=5$) & 530 & 631 & SciPy line search\\
Newton--CG & 228 & 261 (+1869 $Hv$) & linear CG as inner solver\\
\bottomrule
\end{tabular}
\caption{Results for Figure~\ref{fig:rosen100}.}
\end{table}

\begin{figure}[h]
\centering
\includegraphics[width=0.72\textwidth]{ncg_rosen2d.pdf}
\caption{Same comparison in $n=2$, drawn like P18. PR$+$ follows the curved valley in 20 iterations. FR takes many short steps (48 iterations). Steepest descent zig-zags and has not converged after 2000.}
\end{figure}

\paragraph{Observations.} (1)~The qualitative lesson of P18/P20 survives: CG-type directions beat steepest descent by orders of magnitude even on a nonconvex function. (2)~The equivalence of the $\beta$ formulas breaks down: FR and DY are clearly slower than PR$+$/HS/HZ, as the theory predicts. (3)~With more memory (L-BFGS) or curvature information (Newton--CG), iteration counts drop further. This is the trade-off that U20 studies.

%======================================================================
\section{Annotated reading list}
%======================================================================
Marked $\bigstar$ = suggested starting points.

\paragraph{History and foundations.}
\begin{itemize}[leftmargin=1.5em,itemsep=2pt]
\item $\bigstar$ Hestenes \& Stiefel (1952)~\cite{HestenesStiefel1952}: the original paper. It already contains the HS formula and the Krylov optimality.
\item Golub \& O'Leary (1989)~\cite{GolubOLeary1989}: history of CG and Lanczos, 1948--1976.
\item $\bigstar$ Shewchuk (1994)~\cite{Shewchuk1994}: ``CG without the agonizing pain'', a geometric explanation with pictures like P18.
\end{itemize}

\paragraph{Linear CG: theory and finite precision (continues U19).}
\begin{itemize}[leftmargin=1.5em,itemsep=2pt]
\item $\bigstar$ Carson, Liesen \& Strako\v{s} (2024)~\cite{CarsonLiesenStrakos2024}: CG and GMRES through computed examples. Explains why the $\kappa$-bound of P20 can be pessimistic (eigenvalue distribution matters) and how rounding delays convergence.
\item Meurant \& Strako\v{s} (2006)~\cite{MeurantStrakos2006}: Lanczos and CG in finite precision (Acta Numerica).
\item Liesen \& Strako\v{s} (2013)~\cite{LiesenStrakos2013}: book-length treatment of Krylov methods.
\end{itemize}

\paragraph{Preconditioning and computation.}
\begin{itemize}[leftmargin=1.5em,itemsep=2pt]
\item Benzi (2002)~\cite{Benzi2002}: survey of preconditioners. Relevant to the $M$ in \texttt{linear\_cg.py}.
\item Frangella, Tropp \& Udell (2023)~\cite{FrangellaTroppUdell2023}: randomized Nystr\"om preconditioner. The preconditioned condition number becomes constant once the rank is about the effective dimension.
\item Ghysels \& Vanroose (2014)~\cite{GhyselsVanroose2014}: pipelined CG for parallel machines.
\end{itemize}

\paragraph{Nonlinear CG (prepares U20).}
\begin{itemize}[leftmargin=1.5em,itemsep=2pt]
\item $\bigstar$ Hager \& Zhang (2006)~\cite{HagerZhang2006}: \emph{the} survey of nonlinear CG, organized around global convergence.
\item Fletcher \& Reeves (1964)~\cite{FletcherReeves1964}; Polak \& Ribi\`ere (1969)~\cite{PolakRibiere1969}; Polyak (1969)~\cite{Polyak1969}: the founding papers.
\item Al-Baali (1985)~\cite{AlBaali1985}; Powell (1984)~\cite{Powell1984}; Gilbert \& Nocedal (1992)~\cite{GilbertNocedal1992}: the classical convergence results and counterexamples (N\&W §5.2 is based on them).
\item Dai \& Yuan (1999)~\cite{DaiYuan1999}; Hager \& Zhang (2005)~\cite{HagerZhang2005}; Dai \& Kou (2013)~\cite{DaiKou2013}; Hager \& Zhang (2013)~\cite{HagerZhang2013}: modern methods and software.
\item Dai (2011)~\cite{Dai2011}; Narushima \& Yabe (2014)~\cite{NarushimaYabe2014}; Andrei (2020, book)~\cite{Andrei2020}: further surveys.
\item $\bigstar$ Das Gupta, Freund, Sun \& Taylor (2025)~\cite{DasGuptaEtAl2025}: worst-case rates of FR/PRP by computer-assisted proofs. Code is public.
\end{itemize}

\paragraph{CG inside second-order methods (links U7, U15--U17, U21).}
\begin{itemize}[leftmargin=1.5em,itemsep=2pt]
\item $\bigstar$ Steihaug (1983)~\cite{Steihaug1983}: CG with trust regions.
\item Dembo, Eisenstat \& Steihaug (1982)~\cite{DemboEisenstatSteihaug1982}: inexact Newton methods.
\item Nash (2000)~\cite{Nash2000}: survey of truncated Newton methods.
\item Gould, Lucidi, Roma \& Toint (1999)~\cite{GouldLucidiRomaToint1999}: GLTR (Lanczos trust region).
\item Martens (2010)~\cite{Martens2010}: Hessian-free deep learning.
\end{itemize}

\paragraph{Beyond $\R^n$.}
\begin{itemize}[leftmargin=1.5em,itemsep=2pt]
\item Ring \& Wirth (2012)~\cite{RingWirth2012}; Sato \& Iwai (2015)~\cite{SatoIwai2015}; Sato (2022)~\cite{Sato2022}: Riemannian CG.
\end{itemize}

%======================================================================
\section{Conclusion}
%======================================================================
CG generalizes in three ways, and they should be kept separate when reading articles:
\begin{enumerate}[leftmargin=1.5em]
\item \emph{Linear CG on the local quadratic model} (Newton--CG, Steihaug--Toint). All the U18--U19 theory applies to each inner solve. This is the most robust way to use CG for general $f$.
\item \emph{Nonlinear CG}, the same recurrence applied directly to $f$. It is cheap and needs little memory, but finite termination is lost and the various $\beta$'s behave differently. The research questions are global convergence with inexact line searches, sufficient descent, and restarts. The worst-case rate is no better than steepest descent's; the practical advantage is real but asymptotic.
\item \emph{Krylov relatives} for non-SPD linear systems, and CG on manifolds.
\end{enumerate}

%======================================================================
\begin{thebibliography}{99}\small
\bibitem{AlBaali1985} M.~Al-Baali. Descent property and global convergence of the Fletcher--Reeves method with inexact line search. \emph{IMA J. Numer. Anal.}, 5(1):121--124, 1985.
\bibitem{Andrei2020} N.~Andrei. \emph{Nonlinear Conjugate Gradient Methods for Unconstrained Optimization}. Springer, 2020.
\bibitem{Benzi2002} M.~Benzi. Preconditioning techniques for large linear systems: a survey. \emph{J. Comput. Phys.}, 182(2):418--477, 2002.
\bibitem{CarsonLiesenStrakos2024} E.~Carson, J.~Liesen, Z.~Strako\v{s}. Towards understanding CG and GMRES through examples. \emph{Linear Algebra Appl.}, 2024. arXiv:2211.00953.
\bibitem{Dai2011} Y.-H.~Dai. Nonlinear conjugate gradient methods. In \emph{Wiley Encyclopedia of Operations Research and Management Science}, 2011.
\bibitem{DaiKou2013} Y.-H.~Dai, C.-X.~Kou. A nonlinear conjugate gradient algorithm with an optimal property and an improved Wolfe line search. \emph{SIAM J. Optim.}, 23(1):296--320, 2013.
\bibitem{DaiYuan1999} Y.-H.~Dai, Y.~Yuan. A nonlinear conjugate gradient method with a strong global convergence property. \emph{SIAM J. Optim.}, 10(1):177--182, 1999.
\bibitem{Daniel1967} J.~W.~Daniel. The conjugate gradient method for linear and nonlinear operator equations. \emph{SIAM J. Numer. Anal.}, 4:10--26, 1967.
\bibitem{DasGuptaEtAl2025} S.~Das Gupta, R.~M.~Freund, X.~A.~Sun, A.~Taylor. Nonlinear conjugate gradient methods: worst-case convergence rates via computer-assisted analyses. \emph{Math. Program.}, 213:1--49, 2025. doi:10.1007/s10107-024-02127-7.
\bibitem{DemboEisenstatSteihaug1982} R.~S.~Dembo, S.~C.~Eisenstat, T.~Steihaug. Inexact Newton methods. \emph{SIAM J. Numer. Anal.}, 19(2):400--408, 1982.
\bibitem{FaberManteuffel1984} V.~Faber, T.~Manteuffel. Necessary and sufficient conditions for the existence of a conjugate gradient method. \emph{SIAM J. Numer. Anal.}, 21(2):352--362, 1984.
\bibitem{FletcherReeves1964} R.~Fletcher, C.~M.~Reeves. Function minimization by conjugate gradients. \emph{The Computer Journal}, 7(2):149--154, 1964.
\bibitem{FongSaunders2011} D.~C.-L.~Fong, M.~Saunders. LSMR: an iterative algorithm for sparse least-squares problems. \emph{SIAM J. Sci. Comput.}, 33(5):2950--2971, 2011.
\bibitem{FrangellaTroppUdell2023} Z.~Frangella, J.~A.~Tropp, M.~Udell. Randomized Nystr\"om preconditioning. \emph{SIAM J. Matrix Anal. Appl.}, 44(2):718--752, 2023. doi:10.1137/21M1466244.
\bibitem{GhyselsVanroose2014} P.~Ghysels, W.~Vanroose. Hiding global synchronization latency in the preconditioned conjugate gradient algorithm. \emph{Parallel Computing}, 40(7):224--238, 2014.
\bibitem{GilbertNocedal1992} J.~C.~Gilbert, J.~Nocedal. Global convergence properties of conjugate gradient methods for optimization. \emph{SIAM J. Optim.}, 2(1):21--42, 1992.
\bibitem{GolubOLeary1989} G.~H.~Golub, D.~P.~O'Leary. Some history of the conjugate gradient and Lanczos algorithms: 1948--1976. \emph{SIAM Review}, 31(1):50--102, 1989.
\bibitem{GouldLucidiRomaToint1999} N.~I.~M.~Gould, S.~Lucidi, M.~Roma, Ph.~L.~Toint. Solving the trust-region subproblem using the Lanczos method. \emph{SIAM J. Optim.}, 9(2):504--525, 1999.
\bibitem{HagerZhang2005} W.~W.~Hager, H.~Zhang. A new conjugate gradient method with guaranteed descent and an efficient line search. \emph{SIAM J. Optim.}, 16(1):170--192, 2005.
\bibitem{HagerZhang2006} W.~W.~Hager, H.~Zhang. A survey of nonlinear conjugate gradient methods. \emph{Pacific J. Optim.}, 2(1):35--58, 2006.
\bibitem{HagerZhang2013} W.~W.~Hager, H.~Zhang. The limited memory conjugate gradient method. \emph{SIAM J. Optim.}, 23(4):2150--2168, 2013. doi:10.1137/120898097.
\bibitem{HestenesStiefel1952} M.~R.~Hestenes, E.~Stiefel. Methods of conjugate gradients for solving linear systems. \emph{J. Res. Nat. Bur. Standards}, 49(6):409--436, 1952.
\bibitem{LiesenStrakos2013} J.~Liesen, Z.~Strako\v{s}. \emph{Krylov Subspace Methods: Principles and Analysis}. Oxford University Press, 2013.
\bibitem{Martens2010} J.~Martens. Deep learning via Hessian-free optimization. In \emph{Proc.\ ICML}, 2010.
\bibitem{MeurantStrakos2006} G.~Meurant, Z.~Strako\v{s}. The Lanczos and conjugate gradient algorithms in finite precision arithmetic. \emph{Acta Numerica}, 15:471--542, 2006.
\bibitem{NarushimaYabe2014} Y.~Narushima, H.~Yabe. A survey of sufficient descent conjugate gradient methods for unconstrained optimization. \emph{SUT J. Math.}, 50(2):167--203, 2014.
\bibitem{Nash2000} S.~G.~Nash. A survey of truncated-Newton methods. \emph{J. Comput. Appl. Math.}, 124:45--59, 2000.
\bibitem{NocedalWright2006} J.~Nocedal, S.~J.~Wright. \emph{Numerical Optimization}, 2nd ed. Springer, 2006. (Ch.~5 and §7.1.)
\bibitem{PaigeSaunders1975} C.~C.~Paige, M.~A.~Saunders. Solution of sparse indefinite systems of linear equations. \emph{SIAM J. Numer. Anal.}, 12(4):617--629, 1975.
\bibitem{PaigeSaunders1982} C.~C.~Paige, M.~A.~Saunders. LSQR: an algorithm for sparse linear equations and sparse least squares. \emph{ACM Trans. Math. Softw.}, 8(1):43--71, 1982.
\bibitem{PolakRibiere1969} E.~Polak, G.~Ribi\`ere. Note sur la convergence de m\'ethodes de directions conjugu\'ees. \emph{Rev.\ Fran\c{c}aise Informat.\ Recherche Op\'erationnelle}, 3(16):35--43, 1969.
\bibitem{Polyak1969} B.~T.~Polyak. The conjugate gradient method in extremal problems. \emph{USSR Comput. Math. Math. Phys.}, 9(4):94--112, 1969.
\bibitem{Powell1984} M.~J.~D.~Powell. Nonconvex minimization calculations and the conjugate gradient method. In \emph{Numerical Analysis (Dundee 1983)}, LNM 1066, Springer, 122--141, 1984.
\bibitem{RingWirth2012} W.~Ring, B.~Wirth. Optimization methods on Riemannian manifolds and their application to shape space. \emph{SIAM J. Optim.}, 22(2):596--627, 2012.
\bibitem{Saad2003} Y.~Saad. \emph{Iterative Methods for Sparse Linear Systems}, 2nd ed. SIAM, 2003.
\bibitem{Sato2022} H.~Sato. Riemannian conjugate gradient methods: general framework and specific algorithms with convergence analyses. \emph{SIAM J. Optim.}, 32(4), 2022. arXiv:2112.02572.
\bibitem{SatoIwai2015} H.~Sato, T.~Iwai. A new, globally convergent Riemannian conjugate gradient method. \emph{Optimization}, 64(4):1011--1031, 2015.
\bibitem{Shewchuk1994} J.~R.~Shewchuk. An introduction to the conjugate gradient method without the agonizing pain. Tech.\ report, Carnegie Mellon University, 1994.
\bibitem{Steihaug1983} T.~Steihaug. The conjugate gradient method and trust regions in large scale optimization. \emph{SIAM J. Numer. Anal.}, 20(3):626--637, 1983. doi:10.1137/0720042.
\end{thebibliography}

\end{document}
